\documentclass[12pt]{article}
\usepackage[utf8]{inputenc}
\usepackage[T1]{fontenc}
\usepackage{fontspec}
\setmainfont{Times New Roman}
\newfontfamily\hackfont{Hack}
\usepackage{titlesec}
\usepackage{fancyhdr}
\usepackage{listings}
\usepackage{xcolor}
\usepackage{graphicx}
\usepackage{longtable}
\usepackage{caption}
\usepackage{hyperref}
\usepackage{setspace}
\usepackage{geometry}
\geometry{a4paper, margin=1in}
\hypersetup{colorlinks=true, linkcolor=blue, urlcolor=blue}
\renewcommand{\contentsname}{Table of Contents}

\pagestyle{fancy}
\fancyhf{} 							% Header Middle
\rhead{Digital System Design Lab}
\lhead{4bit ALU}
\rfoot{\thepage}

\begin{document}

\begin{titlepage}					% Cover Start
\begin{center}
    
\includegraphics[width=0.35\textwidth]{logo.png} \\
    \vspace{0.5cm}
    \textbf{\Huge \textcolor{teal}{NOTRE DAME} \textcolor{teal}{UNIVERSITY}} \\
    \vspace{0.5cm}
    \textbf{\Huge \textcolor{orange}{BANGLADESH}} \\
    
    \vspace{0.9cm}
    \textbf{\Huge \underline{Digital System Design Lab Task-03}} \\
    
    \vspace{1.5em}
    \begin{flushleft}
    	\textbf{\Large Course Code: CSE-4106} \\
		\vspace{0.3cm}        
        \textbf{\Large Course Title: Digital System Design Lab} \\
        \vspace{0.3cm}
        \textbf{\Large Lab Task Topic: Design 4bit Arithmetic Logic Unit } \\        
    \end{flushleft}
    
    \vspace{1em}
    \begin{flushleft}
        \textbf{\Huge \textcolor{blue}{Submitted by: Group-05}} \\
        \vspace{0.3cm}
        \textbf{\Large Name: Istiak Alam} \\
        \vspace{0.3cm}
        \textbf{\Large ID: 0692230005101005} \\
        \vspace{0.3cm}
        \textbf{\Large Name: Nazifa Tasnim} \\
        \vspace{0.3cm}
        \textbf{\Large ID: 0692230005101007} \\
		\vspace{0.3cm}        
        \textbf{\Large Batch: CSE-20} \\
        \vspace{0.5cm}
        \textbf{\Large Submission Date: }{\Large \textbf{November 9, 2025}\par}
    \end{flushleft}
    \vfill
    \begin{flushleft}
        \textbf{\Huge \textcolor{blue}{Submitted to:}} \\
        \vspace{0.3cm}
        \textbf{\Large Afsana Sharmin Shanta} \\
        \vspace{0.3cm}
        \textbf{\Large Lecturer,} \\
        \vspace{0.3cm}
        \textbf{\Large Notre Dame University Bangladesh}
    \end{flushleft}
\end{center}
\end{titlepage}						% Cover End



\newpage

\section*{Question :}
\subsection*{Design a 4-bit Arithmetic Logic Unit (ALU), which generates the following operations. Also
show four status registers (Sign flag, Carry flag, Overflow flag, Zero flag).}
\begin{center}
\renewcommand{\arraystretch}{1.4}
\begin{tabular}{|l|l|}
\hline
\textbf{Operations} & \textbf{Functions} \\ \hline
$F = B + 1$ & Increment B \\ \hline
$F = B$ & Transfer B \\ \hline
$F = A - B$ & Subtraction \\ \hline
$F = A - B - 1$ & Subtraction with Borrow \\ \hline
$F = A \oplus B$ & Exclusive OR \\ \hline
$F = A \odot B$ & Exclusive NOR \\ \hline
\end{tabular}
\end{center}


\section{Abstract}
This report presents the design and implementation of a 4-bit Arithmetic Logical Unit (ALU) using basic logic gates in Proteus simulation. 
The ALU performs six operations, both arithmetic and logical, based on select inputs and a carry-in signal. 
Additionally, status flags such as Sign Flag (SF), Carry Flag (CF), Overflow Flag (OF), and Zero Flag (ZF) are integrated to indicate result conditions.

\section{Introduction}
An Arithmetic Logical Unit (ALU) is a fundamental component of a CPU that performs arithmetic and logical operations. 
In this project, a 4-bit ALU was designed to perform operations such as increment, transfer, subtraction, XOR, and XNOR based on control signals.

\section{Truth Table}
The ALU uses three select lines ($S_2, S_1, S_0$) and one carry-in input ($Cin$) to select the required operation. 
The complete truth table is shown below:

\begin{center}
\renewcommand{\arraystretch}{1.4}
\begin{tabular}{|c|c|c|c|c|c|l|}
\hline
\textbf{S2} & \textbf{S1} & \textbf{S0} & \textbf{Cin} & \textbf{X} & \textbf{Y} & \textbf{Function} \\ \hline
0 & 0 & 0 & 1 & 0 & B & Increment B $\rightarrow$ B+1 \\ \hline
0 & 0 & 1 & 0 & 0 & B & Transfer B $\rightarrow$ B \\ \hline
0 & 1 & 0 & 1 & A & B' & Subtraction $\rightarrow$ A - B \\ \hline
0 & 1 & 1 & 0 & A & B' & Subtraction with Borrow $\rightarrow$ A - B - 1 \\ \hline
1 & 0 & 0 & x & A & B & Logical XOR $\rightarrow$ A $\oplus$ B \\ \hline
1 & 0 & 1 & x & A & B' & Logical XNOR $\rightarrow$ A $\odot$ B \\ \hline
\end{tabular}
\end{center}

\section{Function Derivations}
The internal signals X, Y, and Z are defined using the select inputs as follows: \\

\begin{itemize}
    \item $X = A \cdot (S_1 \oplus S_2)$
    \item $Y = B \oplus (S_1 + (S_2 \cdot S_0))$
    \item $Z = \overline{S_2} \cdot C_{in}$
\end{itemize}


These functions determine how the ALU inputs are modified for each operation before being processed by the full adder network.

\section{Circuit Diagram}
The 4-bit ALU circuit was designed in Proteus using:
\begin{itemize}
    \item 4-bit Full Adder IC (e.g., 7482)
    \item Logic gates (XOR, OR, NOT, AND)
    \item Select pin for operation selection
    \item Input switches for X, Y, and Cin
    \item \textbf{Number of Logic Gates are 27}
\end{itemize}

\begin{figure}[h!]
    \centering
    \includegraphics[width=0.95\textwidth]{circuit.png}
    \caption{4-bit ALU Circuit Diagram (Proteus Simulation)}
\end{figure}

\section{Status Flags Implementation}
The ALU generates four status flags to indicate the condition of the result:

\begin{itemize}
    \item \textbf{Sign Flag (SF):} Indicates the sign of the result. It is set if the MSB of the result (F3) = 1.
    \[
    SF = F_3
    \]
    
    \item \textbf{Zero Flag (ZF):} Set when all output bits are 0.
    \[
    ZF = \overline{(F_0 + F_1 + F_2 + F_3)}
    \]
    
    \item \textbf{Carry Flag (CF):} Indicates if there is a carry out from the most significant bit.
    \[
    CF = \overline{S_2} \cdot C_{out}
    \]
    
    \item \textbf{Overflow Flag (OF):} Set if there is a signed overflow.
    \[
    OF = (C_3 \oplus C_{out}) \cdot \overline{S_2}
    \]
\end{itemize}

In Proteus, these flags can be displayed using Logicprobe connected to the corresponding outputs of logic circuits:
\begin{itemize}
    \item Use XOR gate for OF
    \item Use NOT and OR gates for ZF
    \item Connect MSB directly for SF
    \item Use the carry output and previous carry from adder for CF
\end{itemize}

\section{Conclusion}
The designed 4-bit ALU successfully performs both arithmetic and logic operations as expected based on select input combinations. 
The inclusion of four status flags allows the ALU to be extended easily for microprocessor-like control applications. 
This experiment provides a clear understanding of how arithmetic logic units function within CPU architecture.

\section{Project File}
\url{https://github.com/Istiaq-Alam/Digital-System-Design-Lab/tree/main/Lab%20Test/Lab%20Test-03}

\end{document}